\section{Git 工作流与 Worktree}

\subsection{Worktree：隔离的工作空间}

\begin{knowledgebox}{Git Worktree 的价值}
Worktree 会在 \texttt{.claude/worktrees/} 下创建一个独立的 Git 分支副本。Claude 在里面怎么折腾都不影响主分支。退出时：没有改动则自动清理，有改动则提示保留或删除。对探索性任务特别有用——不确定某个重构方案能不能行？开个 worktree 试试，不行就扔掉，零风险。
\end{knowledgebox}

\subsection{Git 铁律}

\begin{warningbox}{四条不可违反的 Git 规则}
\begin{enumerate}[leftmargin=1.5em]
  \item \textbf{永远不要让 Claude Code 自动 push}——它可以 commit，但 push 必须由你确认。一旦 push 到远端，回退成本就大了。
  \item \textbf{频繁 commit}——做完一个小功能就 commit。用 \texttt{/rewind} 回退的前提是有 commit 记录。
  \item \textbf{警惕破坏性操作}——\texttt{git reset --hard}、\texttt{git push --force}、\texttt{rm -rf} 没有 undo，可以在权限规则里直接 deny。
  \item \textbf{从 PR 恢复上下文}——自动加载 PR 的改动和讨论作为上下文，适合 code review 或继续别人的工作。
\end{enumerate}
\end{warningbox}

\subsection{本章小结}

Worktree 提供了零风险的实验环境。Git 操作的核心原则是：频繁 commit、禁止自动 push、deny 破坏性命令。

%% ====================================================================
